\documentclass[a4,11pt]{aleph-notas}

% -- Paquetes adicionales
\usepackage{enumitem}
\usepackage{aleph-comandos}


% -- Datos  
\institucion{Facultad de Ciencias Exactas, Naturales y Ambientales}
\carrera{Catálogo STEM}
\asignatura{Álgebra lineal}
\tema{Python no. 1: Matrices y sistemas de ecuaciones}
\autor{Andrés Merino}
\fecha{Periodo 2026-1}

\logouno[0.16\textwidth]{Logos/logoPUCE_04_ac}
\definecolor{colortext}{HTML}{1957d1}
\definecolor{colordef}{HTML}{1957d1}
\fuente{montserrat}


% -- Comandos adicionales
\setlist[enumerate]{label=\roman*.}
\input{formato_ipynb}

\begin{document}

\encabezado


%%%%%%%%%%%%%%%%%%%%%%%%%%%%%%%%%%%%%%
\section{Matrices}
%%%%%%%%%%%%%%%%%%%%%%%%%%%%%%%%%%%%%%


Una manera para trabajar con matrices en Python, es utilizando la librería \texttt{sympy}, con la siguiente línea de comandos:

\begin{pycodigo}
    \begin{ipynbcodigo}\begin{lstlisting}[language=Python]
from sympy import *
    \end{lstlisting}\end{ipynbcodigo}
\end{pycodigo}

Con esto, podemos almacenar y visualizar una matriz de la siguiente manera:
\begin{pycodigo}
    \begin{ipynbcodigo}\begin{lstlisting}[language=Python]
A = Matrix([[1, 2, 3], [4, 5, 6], [-1, 0, 2]])
display(A)
    \end{lstlisting}\end{ipynbcodigo}
    %
    \begin{ipynbsalida}[2mm]
        $\displaystyle \left[\begin{matrix}1 & 2 & 3\\4 & 5 & 6\\-1 & 0 & 2\end{matrix}\right]$
    \end{ipynbsalida}
\end{pycodigo}

    
Cuando se quiera visualizar un único objeto, basta con escribir sunombre al final de la celda:
\begin{pycodigo}
    \begin{ipynbcodigo}\begin{lstlisting}[language=Python]
A
    \end{lstlisting}\end{ipynbcodigo}
    %
    \begin{ipynbsalida}[2mm]
        $\displaystyle \left[\begin{matrix}1 & 2 & 3\\4 & 5 & 6\\-1 & 0 & 2\end{matrix}\right]$
    \end{ipynbsalida}
\end{pycodigo}
    
Para obtener un elemento de la matriz realizamos lo siguiente:
\begin{pycodigo}
    \begin{ipynbcodigo}\begin{lstlisting}[language=Python]
# Entrada de la primera fila y primera columna
A[0, 0]
    \end{lstlisting}\end{ipynbcodigo}
    %
    \begin{ipynbsalida}
        $1$
    \end{ipynbsalida}
    %
    \begin{ipynbcodigo}\begin{lstlisting}[language=Python]
# Entrada de la segunda fila y primera columna
A[1, 0]
    \end{lstlisting}\end{ipynbcodigo}
    %
    \begin{ipynbsalida}
        $4$
    \end{ipynbsalida}
\end{pycodigo}

Para obtener una fila o columna de una matriz realizamos lo siguiente:
\begin{pycodigo}
    \begin{ipynbcodigo}\begin{lstlisting}[language=Python]
# Primera fila
A[0, :]
    \end{lstlisting}\end{ipynbcodigo}
    %
    \begin{ipynbsalida}[2mm]
        $\displaystyle \left[\begin{matrix}1 & 2 & 3\end{matrix}\right]$
    \end{ipynbsalida}
    %
    \begin{ipynbcodigo}\begin{lstlisting}[language=Python]
# Segunda fila
A[1, :]
    \end{lstlisting}\end{ipynbcodigo}
    %
    \begin{ipynbsalida}[2mm]
        $\displaystyle \left[\begin{matrix}4 & 5 & 6\end{matrix}\right]$
    \end{ipynbsalida}
    %
    \begin{ipynbcodigo}\begin{lstlisting}[language=Python]
# Primera columna
A[:, 0]
    \end{lstlisting}\end{ipynbcodigo}
    %
    \begin{ipynbsalida}[2mm]
        $\displaystyle \left[\begin{matrix}1\\4\\-1\end{matrix}\right]$
    \end{ipynbsalida}
    %
    \begin{ipynbcodigo}\begin{lstlisting}[language=Python]
# Segunda columna
A[:, 1]
    \end{lstlisting}\end{ipynbcodigo}
    %
    \begin{ipynbsalida}[2mm]
        $\displaystyle \left[\begin{matrix}2\\5\\0\end{matrix}\right]$
    \end{ipynbsalida}
\end{pycodigo}

Otra forma válida es:
\begin{pycodigo}
    \begin{ipynbcodigo}\begin{lstlisting}[language=Python]
# Primera fila
A.row(0)
    \end{lstlisting}\end{ipynbcodigo}
    %
    \begin{ipynbsalida}[2mm]
        $\displaystyle \left[\begin{matrix}1 & 2 & 3\end{matrix}\right]$
    \end{ipynbsalida}
    %
    \begin{ipynbcodigo}\begin{lstlisting}[language=Python]
# Segunda columna
A.col(1)
    \end{lstlisting}\end{ipynbcodigo}
    %
    \begin{ipynbsalida}[2mm]
        $\displaystyle \left[\begin{matrix}2\\5\\0\end{matrix}\right]$
    \end{ipynbsalida}
\end{pycodigo}

Para comparar la igualdad entre dos matrices, podemos usar el operador \texttt{==}:

\begin{pycodigo}
    \begin{ipynbcodigo}\begin{lstlisting}[language=Python]
# Definimos las matrices
A = Matrix([[1, 0], [0, 1]])
print("A:")
display(A)
B = Matrix([[1, 2], [3, 4]])
print("B:")
display(B)
C = Matrix([[1, 2]])
print("C:")
display(C)
# Comparamos las matrices
print("Es A igual a B: ", A == B)
print("Es B igual a C: ", B == C)
print("Es A igual a A: ", A == A)
    \end{lstlisting}\end{ipynbcodigo}
    %
    \begin{ipynbsalida}
    \begin{Verbatim}
A:
    \end{Verbatim}
    $\displaystyle \left[\begin{matrix}1 & 0\\0 & 1\end{matrix}\right]$
    \begin{Verbatim}
B:
    \end{Verbatim}
    $\displaystyle \left[\begin{matrix}1 & 2\\3 & 4\end{matrix}\right]$
    \begin{Verbatim}
C:
    \end{Verbatim}
    $\displaystyle \left[\begin{matrix}1 & 2\end{matrix}\right]$
    \begin{Verbatim}
Es A igual a B:  False
Es B igual a C:  False
Es A igual a A:  True
    \end{Verbatim}
    \end{ipynbsalida}
\end{pycodigo}

%%%%%%%%%%%%%%%%%%%%%%%%%%%%%%%%%%%%%%
\section{Operaciones de matrices}
%%%%%%%%%%%%%%%%%%%%%%%%%%%%%%%%%%%%%%

Podemos realizar las distintas operaciones de matrices de la siguiente manera:

\begin{pycodigo}
    \begin{ipynbcodigo}\begin{lstlisting}[language=Python]
# Suma de matrices
A+B
    \end{lstlisting}\end{ipynbcodigo}
    %
    \begin{ipynbsalida}[2mm]
        $\displaystyle \left[\begin{matrix}2 & 2\\3 & 5\end{matrix}\right]$
    \end{ipynbsalida}
    %
    \begin{ipynbcodigo}\begin{lstlisting}[language=Python]
# Multiplicación por un escalar
2*B
    \end{lstlisting}\end{ipynbcodigo}
    %
    \begin{ipynbsalida}[2mm]
        $\displaystyle \left[\begin{matrix}2 & 4\\6 & 8\end{matrix}\right]$
    \end{ipynbsalida}
    %
    \begin{ipynbcodigo}\begin{lstlisting}[language=Python]
# Transpuesta
B.T
    \end{lstlisting}\end{ipynbcodigo}
    %
    \begin{ipynbsalida}[2mm]
        $\displaystyle \left[\begin{matrix}1 & 3\\2 & 4\end{matrix}\right]$
    \end{ipynbsalida}
\end{pycodigo}

%%%%%%%%%%%%%%%%%%%%%%%%%%%%%%%%%%%%%%
\subsection{Multiplicación de matrices}

Definamos algunas matrices:
\begin{pycodigo}
    \begin{ipynbcodigo}\begin{lstlisting}[language=Python]
A = Matrix([[1, 2, 3], [3, 2, 1], [1, -1, 0]])
print("A:")
display(A)
B = Matrix([[1, 0, 2], [2, 1, 0], [0, 1, 0]])
print("B:")
display(B)
    \end{lstlisting}\end{ipynbcodigo}
    %
    \begin{ipynbsalida}
    \begin{Verbatim}
A:
    \end{Verbatim}
    $\displaystyle \left[\begin{matrix}1 & 2 & 3\\3 & 2 & 1\\1 & -1 & 0\end{matrix}\right]$
    \begin{Verbatim}
B:
    \end{Verbatim}
    $\displaystyle \left[\begin{matrix}1 & 0 & 2\\2 & 1 & 0\\0 & 1 & 0\end{matrix}\right]$
    \end{ipynbsalida}
\end{pycodigo}

La multiplicación es:
\begin{pycodigo}
    \begin{ipynbcodigo}\begin{lstlisting}[language=Python]
A*B
    \end{lstlisting}\end{ipynbcodigo}
    %
    \begin{ipynbsalida}[2mm]
        $\displaystyle \left[\begin{matrix}5 & 5 & 2\\7 & 3 & 6\\-1 & -1 & 2\end{matrix}\right]$
    \end{ipynbsalida}
\end{pycodigo}


Podemos elevar una matriz a un exponente con el operador \texttt{**}:
\begin{pycodigo}
    \begin{ipynbcodigo}\begin{lstlisting}[language=Python]
print('A^0:')
display(A**0)
print('A^1:')
display(A**1)
print('A^2:')
display(A**2)
    \end{lstlisting}\end{ipynbcodigo}
    %
    \begin{ipynbsalida}[2mm]
    \begin{Verbatim}
A^0:
    \end{Verbatim}

    $\displaystyle \left[\begin{matrix}1 & 0 & 0\\0 & 1 & 0\\0 & 0 & 1\end{matrix}\right]$

    
    \begin{Verbatim}
A^1:
    \end{Verbatim}

    $\displaystyle \left[\begin{matrix}1 & 2 & 3\\3 & 2 & 1\\1 & -1 & 0\end{matrix}\right]$

    
    \begin{Verbatim}
A^2:
    \end{Verbatim}

    $\displaystyle \left[\begin{matrix}10 & 3 & 5\\10 & 9 & 11\\-2 & 0 & 2\end{matrix}\right]$
    \end{ipynbsalida}
\end{pycodigo}
 
\section{Tipos de matrices}\label{tipos-de-matrices}

Podemos obtener las dimensiones de una matriz con el siguiente código:

\begin{pycodigo}
    \begin{ipynbcodigo}\begin{lstlisting}[language=Python]
# Definimos las matrices
A = Matrix([[1, 0, 1], [0, 1, 1]])
print("A:")
display(A)
print("Las dimensiones de A son:", A.shape)
print("El número de filas de A es:", A.shape[0])
print("El número de columnas de A es:", A.shape[1])
    \end{lstlisting}\end{ipynbcodigo}
    %
    \begin{ipynbsalida}
    \begin{Verbatim}
A:
    \end{Verbatim}
    $\displaystyle \left[\begin{matrix}1 & 0 & 1\\0 & 1 & 1\end{matrix}\right]$
    \begin{Verbatim}
Las dimensiones de A son: (2, 3)
El número de filas de A es: 2
El número de columnas de A es: 3
    \end{Verbatim}
    \end{ipynbsalida}
\end{pycodigo}

Ahora, veamos cómo crear algunas matrices especiales:
\begin{itemize}
\item
    Matriz diagonal:
\begin{pycodigo}
    \begin{ipynbcodigo}\begin{lstlisting}[language=Python]
diag(1, 2, 3)
    \end{lstlisting}\end{ipynbcodigo}
    %
    \begin{ipynbsalida}[2mm]
$\displaystyle \left[\begin{matrix}1 & 0 & 0\\0 & 2 & 0\\0 & 0 & 3\end{matrix}\right]$
    \end{ipynbsalida}
\end{pycodigo}

\item
    Matriz identidad:
\begin{pycodigo}
    \begin{ipynbcodigo}\begin{lstlisting}[language=Python]
eye(4)
    \end{lstlisting}\end{ipynbcodigo}
    %
    \begin{ipynbsalida}[2mm]
$\displaystyle \left[\begin{matrix}1 & 0 & 0 & 0\\0 & 1 & 0 & 0\\0 & 0 & 1 & 0\\0 & 0 & 0 & 1\end{matrix}\right]$
    \end{ipynbsalida}
\end{pycodigo}

\item
    Matriz nula:
\begin{pycodigo}
    \begin{ipynbcodigo}\begin{lstlisting}[language=Python]
zeros(2, 4)
    \end{lstlisting}\end{ipynbcodigo}
    %
    \begin{ipynbsalida}[2mm]
$\displaystyle \left[\begin{matrix}0 & 0 & 0 & 0\\0 & 0 & 0 & 0\end{matrix}\right]$
    \end{ipynbsalida}
\end{pycodigo}

\item
    Matriz formada de unos:
\begin{pycodigo}
    \begin{ipynbcodigo}\begin{lstlisting}[language=Python]
ones(4, 2)
    \end{lstlisting}\end{ipynbcodigo}
    %
    \begin{ipynbsalida}[2mm]
$\displaystyle \left[\begin{matrix}1 & 1\\1 & 1\\1 & 1\\1 & 1\end{matrix}\right]$
    \end{ipynbsalida}
\end{pycodigo}
\end{itemize}





\end{document}